% =========================================================================
% Capítulo 4: Cuerpo rígido
% Tema 01: Cuerpo rígido
% Problema 19: Esfera rodante sobre una esfera fija
% Origen: Ayudantía 12 (Sólido Rígido), 2026; corresponde al Ejemplo 5-6 del texto de referencia
% =========================================================================
\subsubsection{Esfera rodante sobre una esfera fija}

Una esfera uniforme de masa $m$ y radio $r$ rueda sin deslizar sobre la superficie de una esfera fija, también de radio $r$, como se muestra en la Figura \ref{fig:c4_e01_p19_esferas}. El ángulo $\phi$ describe la rotación de la esfera móvil respecto a la dirección vertical, mientras que el ángulo $\theta$ describe la posición del centro de masa de la esfera móvil respecto a la vertical. Sabiendo que la gravedad actúa verticalmente hacia abajo:

\begin{enumerate}
    \item Determine la ecuación diferencial de movimiento para la esfera en términos del ángulo $\theta$.
\end{enumerate}

\begin{figure}[!ht]
    \centering
    \begin{tikzpicture}[>=Latex, scale=1.1]
        % Pedestal y tierra de la esfera fija
        \fill[gray!35] (-0.5, -1.2) .. controls (-0.3, -1.5) and (-0.5, -1.8) .. (-0.8, -1.9) -- (0.8, -1.9) .. controls (0.5, -1.8) and (0.3, -1.5) .. (0.5, -1.2) -- cycle;
        \draw[thick] (-0.5, -1.25) .. controls (-0.3, -1.5) and (-0.5, -1.8) .. (-0.8, -1.9);
        \draw[thick] (0.5, -1.25) .. controls (0.3, -1.5) and (0.5, -1.8) .. (0.8, -1.9);
        \fill[pattern=north east lines] (-1.1, -2.15) rectangle (1.1, -1.9);
        \draw[thick] (-1.1, -1.9) -- (1.1, -1.9);

        % Esfera fija (Centro O en (0,0), radio 1.35)
        \shadedraw[inner color=white, outer color=gray!45, draw=black, thick] (0,0) circle (1.35cm);
        \fill (0,0) circle (2pt) node[below left] {$O$};
        \draw[->, thick] (0,0) -- (-35:1.35) node[midway, below left] {$r$};

        % Esfera móvil (Ángulo 40 grados desde la vertical -> 50 grados desde +x)
        \coordinate (C) at (50:2.7);
        \coordinate (P) at (50:1.35);

        \shadedraw[inner color=white, outer color=gray!45, draw=black, thick] (C) circle (1.35cm);
        \fill (C) circle (2pt) node[left] {$C$};
    
        % Radio en esfera móvil y masa m
        \draw[->, thick] (C) -- ($(C)+(-35:1.35)$) node[midway, below left] {$r$};
        \node at ($(C)+(140:0.85)$) {$m$};

        % Línea segmentada O - P - C y vertical de referencia
        \draw[dashed, thick] (0,0) -- (C);
        \draw[dash dot, thick] (0,0) -- (0, 1.5);

        % Arco ángulo theta
        \draw[->, thick] (0, 0.85) arc (90:50:0.85) node[midway, above right] {$\theta$};

        % Punto de contacto P e indicación "Sin deslizar"
        \node[above] at (P) {$P$};
        \draw[thick] ($(P)+(-40:0.15)$) -- ++(-25:1.2) node[right] {Sin deslizar};

        % Arco de rotación phi
        \draw[->, thick] ($(C)+(75:1.65)$) arc (75:30:1.65) node[midway, above right] {$\phi$};

        % Vector gravedad
        \draw[->, thick] (-1.6, 2.6) -- (-1.6, 1.6) node[midway, left] {$\vec{g}$};
    \end{tikzpicture}

    \caption{Esfera de masa $m$ rodando sin deslizar sobre una esfera fija.}
    \label{fig:c4_e01_p19_esferas}
\end{figure}

% --- SOLUCIÓN ---------------------------------------------------------------
\begin{solucion}{Esfera rodante sobre una esfera fija}
\subsubsection*{Cinemática}

Primero, sea $\mathcal{F}$ un marco de referencia fijo a tierra. Elegimos el siguiente sistema de coordenadas fijo en el marco $\mathcal{F}$:
\begin{itemize}
    \item Origen en $O$
    \item $\vec{E}_x =$ A lo largo de $OP$ cuando $\theta = 0$ (verticalmente hacia arriba)
    \item $\vec{E}_z =$ Hacia adentro de la página
    \item $\vec{E}_y = \vec{E}_z \times \vec{E}_x$ (horizontalmente hacia la derecha)
\end{itemize}

A continuación, sea $\mathcal{A}$ un marco de referencia fijo a la dirección de $OC$. Elegimos el siguiente sistema de coordenadas fijo en el marco $\mathcal{A}$:
\begin{itemize}
    \item Origen en $O$
    \item $\vec{e}_r =$ A lo largo de $OP$
    \item $\vec{e}_z = \vec{E}_z =$ Hacia adentro de la página
    \item $\vec{e}_\theta = \vec{e}_z \times \vec{e}_r$
\end{itemize}

La geometría de las bases $\{\vec{E}_x, \vec{E}_y, \vec{E}_z\}$ y $\{\vec{e}_r, \vec{e}_\theta, \vec{e}_z\}$ se muestra a continuación:

\begin{center}
\begin{tikzpicture}[scale=1.1, >=Stealth]
    \draw[->, thick] (0,0) -- (0,2.2) node[left] {$\vec{E}_x$};
    \draw[->, thick] (0,0) -- (2.2,0) node[below] {$\vec{E}_y$};
    \draw[->, thick] (0,0) -- (50:2.2) node[right] {$\vec{e}_r$};
    \draw[->, thick] (0,0) -- (-40:2.2) node[below left] {$\vec{e}_\theta$};

    % Símbolo hacia adentro de la página
    \filldraw[fill=white, thick] (0,0) circle (4pt);
    \draw[thick] (-2.8pt,-2.8pt) -- (2.8pt,2.8pt);
    \draw[thick] (-2.8pt,2.8pt) -- (2.8pt,-2.8pt);
    \node[left] at (-0.15,0) {$\vec{e}_z, \vec{E}_z$};

    % Arcos de ángulo theta
    \draw[->, thin] (0,1.3) arc (90:50:1.3);
    \node at (70:1.5) {$\theta$};
    \draw[->, thin] (1.3,0) arc (0:-40:1.3);
    \node at (-20:1.5) {$\theta$};
\end{tikzpicture}
\end{center}

A partir de la geometría de las bases, tenemos:
\begin{align}
    \vec{E}_x &= \cos\theta\,\vec{e}_r - \sin\theta\,\vec{e}_\theta \\
    \vec{E}_y &= \sin\theta\,\vec{e}_r + \cos\theta\,\vec{e}_\theta
\end{align}

Finalmente, sea $\mathcal{R}$ el marco de referencia de la esfera móvil (no es necesario especificar un sistema de coordenadas fijo en $\mathcal{R}$). La posición del centro de masa de la esfera móvil es:
\begin{equation}
    \bar{\vec{r}} = 2r\vec{e}_r
\end{equation}
Además, la velocidad angular del marco $\mathcal{A}$ en $\mathcal{F}$ y la velocidad angular del marco $\mathcal{R}$ en $\mathcal{F}$ (dado que $\phi$ se mide respecto a la vertical fija) están dadas por:
\begin{align}
    {}^{\mathcal{F}}\vec{\omega}^{\mathcal{A}} &= \dot{\theta}\vec{e}_z \\
    {}^{\mathcal{F}}\vec{\omega}^{\mathcal{R}} &= \dot{\phi}\vec{e}_z
\end{align}

Aplicando el teorema del transporte entre los marcos $\mathcal{A}$ y $\mathcal{F}$, la velocidad del centro de masa de la esfera es:
\begin{equation}
    {}^{\mathcal{F}}\bar{\vec{v}} = \frac{{}^{\mathcal{F}}d\bar{\vec{r}}}{dt} = \frac{{}^{\mathcal{A}}d\bar{\vec{r}}}{dt} + {}^{\mathcal{F}}\vec{\omega}^{\mathcal{A}} \times \bar{\vec{r}} = \vec{0} + \dot{\theta}\vec{e}_z \times (2r\vec{e}_r) = 2r\dot{\theta}\vec{e}_\theta
\end{equation}

Aplicando nuevamente el teorema del transporte a ${}^{\mathcal{F}}\bar{\vec{v}}$ entre $\mathcal{A}$ y $\mathcal{F}$, obtenemos la aceleración del centro de masa:
\begin{equation}
    {}^{\mathcal{F}}\bar{\vec{a}} = \frac{{}^{\mathcal{F}}d}{dt}({}^{\mathcal{F}}\bar{\vec{v}}) = \frac{{}^{\mathcal{A}}d}{dt}({}^{\mathcal{F}}\bar{\vec{v}}) + {}^{\mathcal{F}}\vec{\omega}^{\mathcal{A}} \times {}^{\mathcal{F}}\bar{\vec{v}} = 2r\ddot{\theta}\vec{e}_\theta + \dot{\theta}\vec{e}_z \times (2r\dot{\theta}\vec{e}_\theta) = -2r\dot{\theta}^2\vec{e}_r + 2r\ddot{\theta}\vec{e}_\theta
\end{equation}

Una segunda expresión para ${}^{\mathcal{F}}\bar{\vec{v}}$ se obtiene observando que los puntos $C$ y $P$ pertenecen a la esfera móvil $\mathcal{R}$:
\begin{equation}
    {}^{\mathcal{F}}\bar{\vec{v}} = {}^{\mathcal{F}}\vec{v}_P^{\mathcal{R}} + {}^{\mathcal{F}}\vec{\omega}^{\mathcal{R}} \times (\bar{\vec{r}} - \vec{r}_P)
\end{equation}
donde $\bar{\vec{r}} - \vec{r}_P = r\vec{e}_r$. Por la condición de rodadura sin deslizamiento, ${}^{\mathcal{F}}\vec{v}_P^{\mathcal{R}} = \vec{0}$. En consecuencia:
\begin{equation}
    {}^{\mathcal{F}}\bar{\vec{v}} = \dot{\phi}\vec{e}_z \times r\vec{e}_r = r\dot{\phi}\vec{e}_\theta
\end{equation}
Igualando ambas expresiones de ${}^{\mathcal{F}}\bar{\vec{v}}$:
\begin{equation}
    2r\dot{\theta} = r\dot{\phi} \implies \dot{\phi} = 2\dot{\theta} \implies \ddot{\phi} = 2\ddot{\theta}
\end{equation}
Por lo tanto, la velocidad angular de la esfera superior en el marco $\mathcal{F}$ es:
\begin{equation}
    {}^{\mathcal{F}}\vec{\omega}^{\mathcal{R}} = 2\dot{\theta}\vec{e}_z
\end{equation}

\subsubsection*{Cinética}

Resolveremos el problema aplicando ambas leyes de Euler utilizando el centro de masa de la esfera como punto de referencia. El diagrama de cuerpo libre (DCL) se muestra a continuación:

\begin{center}
\begin{tikzpicture}[scale=1.0, >=Stealth]
    \def\R{1.35}
    \shadedraw[ball color=gray!15, shading=ball, opacity=0.85, thick] (0,0) circle (\R cm);
    \fill (0,0) circle (2pt);
    % Peso
    \draw[->, very thick] (0,0) -- (0,-1.0) node[midway, left] {$m\vec{g}$};
    % Punto de contacto P en angulo -130 grados
    \coordinate (P) at (-130:\R);
    % Fuerza Normal N (en direccion e_r, angulo 50)
    \draw[<-, very thick] (P) -- ++(-130:1.1) node[left] {$\vec{N}$};
    % Fuerza de roce F_f (en direccion e_theta, angulo -40)
    \draw[->, very thick] (P) -- ++(-40:1.1) node[below] {$\vec{F}_f$};
\end{tikzpicture}
\end{center}

Expresando la fuerza normal $\vec{N}$, la fuerza de rodadura (roce) $\vec{F}_f$ y la fuerza de gravedad $m\vec{g}$ en la base $\{\vec{e}_r, \vec{e}_\theta, \vec{e}_z\}$:
\begin{align}
    \vec{N} &= N\vec{e}_r \\
    \vec{F}_f &= F_f\vec{e}_\theta \\
    m\vec{g} &= -mg\vec{E}_x = -mg(\cos\theta\,\vec{e}_r - \sin\theta\,\vec{e}_\theta) = -mg\cos\theta\,\vec{e}_r + mg\sin\theta\,\vec{e}_\theta
\end{align}

\textit{1. Aplicación de la 1ª Ley de Euler al centro de masa:}
La fuerza resultante $\vec{F}$ sobre la esfera es:
\begin{equation}
    \vec{F} = (N - mg\cos\theta)\vec{e}_r + (F_f + mg\sin\theta)\vec{e}_\theta
\end{equation}
Igualando $\vec{F} = m\,{}^{\mathcal{F}}\bar{\vec{a}}$:
\begin{equation}
    (N - mg\cos\theta)\vec{e}_r + (F_f + mg\sin\theta)\vec{e}_\theta = -2mr\dot{\theta}^2\vec{e}_r + 2mr\ddot{\theta}\vec{e}_\theta
\end{equation}
Lo cual entrega dos ecuaciones escalares:
\begin{align}
    N - mg\cos\theta &= -2mr\dot{\theta}^2 \\
    F_f + mg\sin\theta &= 2mr\ddot{\theta} \label{eq:c4_e01_p19_f}
\end{align}

\textit{2. Aplicación de la 2ª Ley de Euler respecto al centro de masa:}
Como $\vec{N}$ y $m\vec{g}$ pasan por el centro de masa, solo $\vec{F}_f$ produce momento respecto a $C$. Como $\vec{r}_P - \bar{\vec{r}} = -r\vec{e}_r$:
\begin{equation}
    \bar{\vec{M}} = (\vec{r}_P - \bar{\vec{r}}) \times \vec{F}_f = -r\vec{e}_r \times F_f\vec{e}_\theta = -rF_f\vec{e}_z
\end{equation}
El momento angular respecto al centro de masa es ${}^{\mathcal{F}}\bar{\vec{H}} = \bar{\vec{I}}^{\mathcal{R}} \cdot {}^{\mathcal{F}}\vec{\omega}^{\mathcal{R}}$. Dado que para una esfera uniforme $\bar{I}_{zz} = \frac{2}{5}mr^2$:
\begin{equation}
    {}^{\mathcal{F}}\bar{\vec{H}} = \bar{I}_{zz}(2\dot{\theta}\vec{e}_z) = \frac{2}{5}mr^2(2\dot{\theta}\vec{e}_z) = \frac{4}{5}mr^2\dot{\theta}\vec{e}_z
\end{equation}
Derivando en el marco $\mathcal{F}$:
\begin{equation}
    \frac{{}^{\mathcal{F}}d}{dt}({}^{\mathcal{F}}\bar{\vec{H}}) = \frac{4}{5}mr^2\ddot{\theta}\vec{e}_z
\end{equation}
Aplicando $\bar{\vec{M}} = \frac{{}^{\mathcal{F}}d}{dt}({}^{\mathcal{F}}\bar{\vec{H}})$:
\begin{equation}
    -rF_f = \frac{4}{5}mr^2\ddot{\theta} \implies F_f = -\frac{4}{5}mr\ddot{\theta}
\end{equation}

\subsubsection*{Ecuación Diferencial de Movimiento}

Sustituyendo $F_f$ en la ecuación (\ref{eq:c4_e01_p19_f}):
\begin{equation}
    -\frac{4}{5}mr\ddot{\theta} + mg\sin\theta = 2mr\ddot{\theta} \implies \frac{14}{5}mr\ddot{\theta} = mg\sin\theta
\end{equation}
Simplificando, obtenemos la ecuación diferencial de movimiento para la esfera:
\respuesta{$\displaystyle \ddot{\theta} - \frac{5g}{14r}\sin\theta = 0$}
\end{solucion}
